\subsection{Eliminating a split characteristic-polynomial fibre}
\label{tr:parabolic-descent}
Here is a local proof of assertion (ii) of \S\ref{tr:geometry}.
Let $F$ be the local field at the chosen split place, $K=\mathrm{GL}_r(\mathcal O_F)$,
and let $R=R_1R_2$ have coprime monic factors of positive degrees $a,b$
and nonzero constant terms. Write $G=\mathrm{GL}_r(F)$, $M=\mathrm{GL}_a(F)\times
\mathrm{GL}_b(F)$ and $U$ for the upper block radical, identified with
the additive space of $a$ by $b$ matrices. Suppose the compactly supported
locally constant test $f_0$ has
\[
 \int_U f_0(k m u k^{-1})\,du=0\qquad (k\in K,\ m\in M).
\]
For the inflated character idempotent this follows from the radical sums
of \S\ref{tr:label}. Indeed $mu\in K$ forces both diagonal blocks to be
integral and invertible, and then $u\in U(\mathcal O_F)$.
Reduction integrates each residue class with the same measure.
The same conclusion holds for any fixed left and right translates of the
test if the finite radical sums vanish for arbitrary left and right
multipliers: by $G=K M U= M U K$ the outside factors reduce to compact
multipliers and an affine change of the variable in $U$.

\begin{lemma}
Every conjugation-invariant distribution on $G$ supported on the fibre of
$R$ annihilates $f_0$.
\end{lemma}
\begin{proof}
On this closed fibre, the ordered complementary subspaces
$\ker R_1(g),\ker R_2(g)$ have dimensions $a,b$. The projectors onto them
are fixed polynomials in $g$ by B\'ezout's identity. The map to complementary
pairs is therefore continuous and equivariant; the space of pairs is
$G/M$. Over the coordinate pair the fibre is the space $X$ of diagonal
blocks $(A,B)$ with their indicated characteristic polynomials. A test
vanishing on the closed fibre is annihilated by the distribution, since
its compact open support misses that fibre. Tests on the fibre extend
locally constantly with compact support, by a finite compact open cover.
Thus Frobenius descent for locally compact totally disconnected spaces
applies to the distribution on this fibre: $G$ and $M$ are both
unimodular. See Aizenbud, Gourevitch, Rallis and Schiffmann,
\emph{Multiplicity one theorems}, Theorem 2.2~\cite{agrs}.

Concretely the descended distribution on $X$ acts on the test
\[
 (A,B)\longmapsto \int_K\int_U
 f_0(k u\operatorname{diag}(A,B)u^{-1}k^{-1})\,du\,dk .
\]
This is the quotient integral in descent: in the Iwasawa coordinates
$g=kum$ one disintegrates Haar measure as $dk\,du\,dm$, up to a fixed
positive scalar, and absorbs $m$ by invariance of the distribution
on $X$. The displayed integral is compactly supported on $X$.
Indeed the compact support bounds the diagonal and off-diagonal blocks
after conjugation by $K$; the inverse of the map $Z\mapsto ZB-AZ$ is
uniformly bounded there. Its determinant, up to sign, is the nonzero
constant $\operatorname{Res}(R_1,R_2)$, as seen by triangularizing
$A,B$ separately over a splitting field. The same bounds make the
integral locally constant and justify testing the distribution under
the integral.

For fixed diagonal blocks, the upper-right entry of the conjugate by
$u(Z)$ is $ZB-AZ$. Changing variables from $Z$ to $A^{-1}(ZB-AZ)$
expresses the inner integral as a constant times the radical integral
in the hypothesis, hence it vanishes. This argument also covers
nonsemisimple diagonal blocks.
\end{proof}
